\documentclass[11pt]{article}
\usepackage{setspace}
\usepackage{geometry}
\geometry{margin=1in}
\setstretch{1.15}

\begin{document}

% --- Applicant Information ---
\noindent
Zhongli Wang \\
Phone: +49 160 9541 6015 \\
Email: nkwangzhongli@gmail.com \\
Website: https://sites.google.com/view/zhongli-wang \\
Bielefeld, Germany \\

\vspace{0.5cm}

% --- Recipient Information ---
\noindent
Ms.\ Edda B\"ohning \\
SOCIUM -- Research Center on Inequality and Social Policy \\
University of Bremen \\
P.O.\ Box 33 04 40 \\
28334 Bremen \\
Germany \\

\vspace{0.5cm}

% --- Subject Line ---
\noindent
\textbf{Application for the position of Wissenschaftliche:r Mitarbeiter:in (Postdoc/Doktorand:in) -- Policy Evaluation / SFB 1342} \\

\vspace{0.3cm}

\noindent Dear Ms.\ B\"ohning, \\
\noindent 
Dear Prof.\ Fehrler,\\[0.2cm]

I am writing to apply for the position of wissenschaftliche:r Mitarbeiter:in (Postdoc/Doktorand:in) affiliated with the working group ``Economics of Social Policy'' (Prof.\ Sebastian Fehrler), the Chair of Macroeconomics (Prof.\ Torben Klarl), and the Collaborative Research Center SFB 1342 ``Global Dynamics of Social Policy.'' 
I am currently completing my Ph.D.\ in Economics at Bielefeld University and Universit\`a Cattolica del Sacro Cuore within the Marie Sk{\l}odowska--Curie EPOC programme, and I plan to defend before the start of employment. 
My research combines applied microeconometrics and computational modeling to study how policies and institutional constraints shape individual behavior and market outcomes. 
With strong quantitative skills, experience in building reproducible empirical pipelines, and C1-level German and English proficiency, I am motivated to contribute to SOCIUM's research agenda and to the empirical work of SFB 1342.

My doctoral research focuses on decision-making under frictions and heterogeneity, using a mix of empirical analysis, simulation-based estimation, and computational methods. 
Substantively, I study how incentives and constraints affect behavior and welfare, and how policy interventions can be evaluated through counterfactual analysis. 
Methodologically, I have extensive experience with micro data handling, model-based and reduced-form evaluation strategies, and high-performance computational workflows in Python (with complementary use of R/Stata as needed). 
% This background provides a strong foundation for policy evaluation with a micro- and macroeconomic focus, including the careful construction of datasets, transparent identification arguments, and robustness analysis.
In parallel with my dissertation, I have developed several side projects that strengthen my fit for an empirically oriented, collaborative research environment. I coauthored \textit{Polytope Fraud Theory} (IRFA, 2025), which builds and validates machine-learning based indicators from firm financial data, sharpening my skills in data construction, prediction, and interpretability. I also coauthored \textit{Agentic Workflows for Economic Research}, which designs reproducible AI-assisted pipelines for literature review, data extraction, and analysis, reinforcing my commitment to transparent research practices. 
% In addition, I have been involved in organizing research workshops and interdisciplinary reading groups, experiences that translate well to seminar activities, project coordination, and research communication within SOCIUM and the SFB environment.

The advertised position is particularly appealing because it connects rigorous empirical policy evaluation to broader questions of social policy dynamics. 
I am eager to contribute to work in the ``Economics of Social Policy'' group on the design and impact analysis of social policy programmes, including settings in the Global South where measurement challenges, heterogeneous treatment effects, and implementation constraints make credible evaluation both difficult and impactful. 
I am also interested in the group's second research area on strategic decision-making and group behavior. 
While my prior work has primarily used computational approaches, I am keen to strengthen the experimental and game-theoretic side of my profile by working on projects that combine theory, lab/survey evidence, and empirical evaluation.

In addition, I welcome the connection to macroeconomic perspectives through Prof.\ Klarl's Chair and the SFB 1342 environment. 
My research naturally interfaces with macro questions when micro-level behavioral responses and institutional settings aggregate into system-level outcomes. 
I would be excited to develop projects that link micro evidence on program design and incentives to macro-relevant outcomes such as labor supply, human capital accumulation, and distributional dynamics, and to contribute to SFB-related collaborative work with international partners, particularly the University of Graz. 
I am fully open to research stays in Graz and value sustained collaboration in interdisciplinary teams.

Alongside research, I am prepared to take on the teaching responsibilities (3.5 SWS) in Bachelor and Master programmes, including instruction in scientific methods and quantitative tools. 
My teaching experience includes two years as a teaching assistant at the University of Zurich and the development of technology-supported teaching components (including automated feedback systems) for large courses. 
In my teaching, I emphasize clarity, research-based learning, and hands-on training in empirical methods, so that students learn to formulate questions, work with data responsibly, and communicate results with precision.
More broadly, the University of Bremen's emphasis on interdisciplinarity, research-based learning, and openness to unconventional approaches resonates strongly with my working style. 
I would welcome the opportunity to contribute to the SOCIUM community through seminars, workshops, and collaborative projects, and to help strengthen the empirical and computational capacity of ongoing research activities, including data management and reproducible research practices.

Thank you so much for considering my application! 
I would be delighted to discuss how my quantitative background and research interests can contribute to the working groups, the SFB 1342 research programme, and the teaching mission at the University of Bremen.

\vspace{6mm}

\noindent Sincerely,\\[12mm]
\noindent Zhongli Wang

\end{document}